\chapter{Теоретические основы сверточных нейронных сетей}

\section{Математическая модель двумерной дискретной свертки}

Классические полносвязные нейронные сети (Multi-Layer Perceptron, MLP) при обработке визуальных данных сталкиваются с проблемой экспоненциального роста числа обучаемых параметров и разрушения двумерной пространственной топологии изображения. Вытягивание матрицы пикселей в одномерный вектор утрачивает информацию о взаимном расположении соседних точек.

Сверточные нейронные сети (Convolutional Neural Networks, CNN) преодолевают указанное ограничение за счет двух ключевых архитектурных принципов: \textit{локальности рецептивных полей} и \textit{разделения весов} (weight sharing).

Операция двумерной дискретной свертки входной матрицы изображения $I \in \mathbb{R}^{H \times W}$ с обучаемым ядром (фильтром) $K \in \mathbb{R}^{k_h \times k_w}$ определяется выражением:
\begin{equation}
    S(i, j) = (I * K)(i, j) = \sum_{m=0}^{k_h - 1} \sum_{n=0}^{k_w - 1} I(i + m, j + n) K(m, n) + b,
    \label{eq:convolution}
\end{equation}
где $b \in \mathbb{R}$ --- смещение (bias), $S(i, j)$ --- результирующее значение в карте признаков (feature map).

При наличии $C_{\text{in}}$ входных каналов (для цветного RGB-изображения $C_{\text{in}} = 3$) и набора из $C_{\text{out}}$ фильтров, сверточный слой преобразует тензор размера $(H, W, C_{\text{in}})$ в тензор $(H', W', C_{\text{out}})$. Каждое ядро специализируется на выявлении определенных локальных паттернов: границы, градиенты яркости на первых слоях, текстуры и геометрические формы на промежуточных, целостные структуры лепестков и соцветий на глубоких уровнях абстракции.

После операции свертки к каждому элементу карты признаков применяется поэлементная нелинейная функция активации. В данной работе используется выпрямленная линейная единица $\text{ReLU}$ (Rectified Linear Unit):
\begin{equation}
    f(x) = \max(0, x),
    \label{eq:relu}
\end{equation}
производная которой равна единице при $x > 0$ и нулю при $x \le 0$. Данная функция исключает эффект насыщения и затухания градиентов, характерный для сигмоиды и гиперболического тангенса.

Для снижения размерности карт признаков и формирования инвариантности к малым геометрическим сдвигам применяется операция максимального объединения (MaxPooling):
\begin{equation}
    P(i, j) = \max_{0 \le m < p, \, 0 \le n < p} S(i \cdot s + m, \, j \cdot s + n),
    \label{eq:maxpool}
\end{equation}
где $p \times p$ --- размер окна пулинга (в работе $2 \times 2$), $s$ --- шаг (stride, $s = 2$). Пулинг сокращает пространственное разрешение в два раза по каждому измерению, уменьшая вычислительную сложность последующих слоев.

\section{Механизм пакетной нормализации (BatchNormalization)}

При последовательном прохождении сигнала через глубокие слои нейросети распределение входов каждого слоя непрерывно изменяется по мере обновления весов предыдущих слоев. Данный феномен известен как \textit{внутренний ковариационный сдвиг} (Internal Covariate Shift). Слоям приходится адаптироваться к изменяющимся масштабам и смещениям, что вынуждает использовать предельно малые скорости обучения и замедляет сходимость.

Слой пакетной нормализации (BatchNormalization, предложенный С.~Иоффе и К.~Сегеди) стабилизирует процесс оптимизации. Для каждого мини-батча $\mathcal{B} = \{x_1, x_2, \dots, x_m\}$ размера $m$ вычисляются выборочные среднее и дисперсия:
\begin{equation}
    \mu_\mathcal{B} = \frac{1}{m} \sum_{i=1}^m x_i, \qquad \sigma_\mathcal{B}^2 = \frac{1}{m} \sum_{i=1}^m (x_i - \mu_\mathcal{B})^2.
    \label{eq:bn_mean_var}
\end{equation}

Далее выполняется нормирование с добавлением малой константы $\epsilon = 10^{-3}$ для предотвращения деления на ноль:
\begin{equation}
    \hat{x}_i = \frac{x_i - \mu_\mathcal{B}}{\sqrt{\sigma_\mathcal{B}^2 + \epsilon}}.
    \label{eq:bn_norm}
\end{equation}

Для сохранения выразительной способности сети вводятся два обучаемых параметра --- масштаб $\gamma$ и сдвиг $\beta$:
\begin{equation}
    y_i = \gamma \hat{x}_i + \beta.
    \label{eq:bn_scale_shift}
\end{equation}

В сверточных слоях нормализация производится по пространственным осям и мини-батчу одновременно, сохраняя единые масштабирующие коэффициенты для каждой карты признаков. На этапе инференса используются накопленные экспоненциальные скользящие средние (moving averages), что обеспечивает детерминированность вычислений.

Пакетная нормализация выполняет две важнейшие функции:
\begin{enumerate}
    \item \textbf{Сглаживание ландшафта оптимизации:} функция потерь становится более липшицевой, градиенты не взрываются и не затухают, позволяя безопасно повышать скорость обучения.
    \item \textbf{Слабая регуляризация:} стохастические флуктуации статистик по мини-батчам вносят полезный шум, препятствующий сонастройке весов и снижающий эффект переобучения без необходимости принудительного обнуления нейронов, как в Dropout.
\end{enumerate}

\section{Оптимизация стохастическим градиентным спуском с моментом}

Стандартный градиентный спуск обновляет веса $\theta \in \mathbb{R}^d$ вдоль антиградиента эмпирического риска по всей выборке: $\theta_{t} = \theta_{t-1} - \eta \nabla \mathcal{L}(\theta_{t-1})$. В задачах глубокого обучения вычисление градиента по всей выборке ресурсоемко, поэтому применяется стохастический градиентный спуск (SGD), где градиент оценивается по случайному мини-батчу.

Однако в сложных невыпуклых овражистых поверхностях простой SGD испытывает сильные паразитные колебания поперек крутых стенок и медленно продвигается вдоль пологого дна оврага к минимуму. Для подавления этих колебаний используется классический импульс (Momentum) и модификация Нестерова (NAG):
\begin{align}
    v_t &= \beta v_{t-1} + \eta \nabla_\theta \mathcal{L}(\theta_{t-1} - \beta v_{t-1}), \label{eq:nesterov_v} \\
    \theta_t &= \theta_{t-1} - v_t, \label{eq:nesterov_theta}
\end{align}
где $\beta \in [0, 1)$ --- коэффициент импульса (в работе $\beta = 0.9$), $v_t$ --- вектор накопленной скорости.

Поправка Нестерова вычисляет градиент не в текущей точке $\theta_{t-1}$, а в «предсказанной» точке шага по инерции $\theta_{t-1} - \beta v_{t-1}$, что действует как упреждающее адаптивное демпфирование и существенно сокращает перелеты через зоны оптимума.
